\chapter{2000 Nobel Prize in Economics}
\textbf{Laureates: }James Heckman (USA), Daniel McFadden (USA)

\section*{Reason for Award}
For his development of theory and methods for analyzing selective samples, and
for Daniel McFadden's development of theory and methods for analyzing discrete
choice

\section{What They Did}
James Heckman and Daniel McFadden were honored for pioneering contributions to
the econometric analysis of microdata that address two fundamental and pervasive
problems in empirical research: selection bias and discrete decision-making.

Heckman developed what has become known as the "Heckman correction" (also
called the inverse Mills ratio method) to address sample selection bias. The
problem arises when researchers observe outcomes only for a non-randomly
selected subsample of the population. A classic example is the study of wages:
researchers typically observe wages only for individuals who are employed, not
for those who have chosen not to participate in the labor force. When the
decision to participate is correlated with unobserved factors that also affect
wages, standard ordinary least squares regression produces biased and
inconsistent estimates.

Heckman showed how selection can be modeled through a two-stage procedure. First,
a probit selection equation estimates the probability of being included in the
sample. Second, an estimated inverse Mills ratio (also called the Heckman lambda)
is included as an additional regressor in the outcome equation. Under assumptions
such as a correctly specified selection equation, an appropriate exclusion
restriction or functional form, and a suitable error structure, this accounts for
selection-induced correlation and can produce consistent estimates. It is not a
general correction that works without those assumptions.

Beyond sample selection, Heckman made important contributions to the evaluation
of social programs and to the analysis of heterogeneous treatment effects. He
developed methods for learning from non-experimental data, recognizing that
randomized controlled trials are not always feasible or ethical. His work on
life-cycle models of skill formation emphasized self-productivity and dynamic
complementarity: early investments can raise the returns to later investments,
although the size and timing of returns depend on the setting.

McFadden developed the econometric theory of discrete-choice models. When
individuals choose among a finite set of alternatives--such as modes of
transportation, employers, neighborhoods, or consumer products--their choices
can be modeled using random-utility theory. McFadden showed that the multinomial
logit model arises from assuming that the unobserved utility components are
independently and identically distributed according to a Type I extreme-value
distribution. This assumption also implies the restrictive independence of
irrelevant alternatives property.

This result made discrete-choice modeling computationally tractable, although the
basic logit model remains restrictive. McFadden established important statistical
properties of maximum-likelihood estimation in this context, including conditions
for consistency and asymptotic normality. The conditional-logit model extended
these ideas to alternatives described by alternative-specific attributes. Later
mixed-logit models accommodated richer substitution patterns and unobserved
heterogeneity. McFadden's 1974 conditional-logit paper and subsequent work became
foundational references in the field.

\section{Impact on Economic Thinking}
The Heckman correction gave researchers a formal way to address a persistent and
often invisible source of bias in labor economics, health economics, education
research, and program evaluation. Studies of wage determination, labor-force
participation, and social programs can be misleading when individuals self-select
into the observed outcome; the correction helps when its identifying assumptions
are credible.

McFadden's discrete-choice models became foundational in empirical work on
consumer behavior, transportation planning, marketing research, and
environmental valuation. The logit model is computationally tractable, while
its assumptions must be assessed against the substitution patterns and
heterogeneity in the application. McFadden's work provided statistical machinery
for estimating choice models under explicit behavioral and distributional
assumptions.

Together, their methods enabled researchers to evaluate policy programs, study
discrimination, and forecast demand with more defensible empirical models. Their
work belongs to the broader development of identification and causal inference;
instrumental variables, regression discontinuity, and difference-in-differences
also have distinct intellectual origins and should not be presented as direct
extensions of Heckman's methods.

Heckman's life-cycle models of skill formation, developed with collaborators,
emphasized that early investments can affect the productivity of later investments
and that skill formation is dynamic. This insight informed education policy,
including debates about preschool and later remedial programs, but the relative
returns depend on program quality, age, population, and broader economic context.

\section{Modern Perspectives Today}
Heckman and McFadden's methodologies are widely used across the social sciences.
In labor economics, methods for selection and treatment-effect analysis remain
important for evaluating employment-training programs, labor-market policies, and
the returns to education.

Discrete-choice models inform transportation policy, housing-market analysis,
environmental valuation, and marketing strategy. They are used to estimate
willingness to pay for policy changes, forecast travel demand under different
infrastructure scenarios, and study consumer and firm decisions. Health
economists use discrete-choice experiments to elicit patient preferences for
treatment options and healthcare attributes.

Machine-learning methods are sometimes combined with selection, treatment-effect,
and discrete-choice models in high-dimensional settings. Modern causal-inference
methods, including causal forests and double machine learning, have their own
literatures but share the broader goal of making identification and heterogeneity
explicit. Their combined legacy remains an important part of empirical
microeconomics.

\subsection*{Key References}
\begin{itemize}
    \item Heckman, J. J. (1979). ``Sample Selection Bias as a Specification Error.'' \textit{Econometrica}, 47(1), 153--161.
    \item McFadden, D. (1974). ``Conditional Logit Analysis of Qualitative Choice Behavior.'' In P. Zarembka (Ed.), \textit{Frontiers in Econometrics}, 105--142. Academic Press.
\end{itemize}
